\section{Grothendieck universes}\label{sec:grothendieck_universes}

Instead of having one single universe, we can have multiple universes where each is contained in another one. The upside of this is that we can do category theory formally within set theory --- see the discussions in \cref{def:small_category}. The downside of this is that, unlike models of \hyperref[def:zfc]{\logic{ZFC}}, models of \hyperref[def:zfcu]{\logic{ZFC+U}} (\logic{ZFC} with the axiom that every set is contained in some \hyperref[def:grothendieck_universe]{Grothendieck universe}) are much less studied. In particular, this axiom requires the existence of an unbounded hierarchy of \hyperref[rem:strongly_inaccessible_cardinal]{regular strong limit cardinal}, unlike \( \logic{ZFC} \) for which only one such cardinal is sufficient.

\begin{definition}\label{def:grothendieck_universe}\mcite[22]{MacLane1998CategoryTheory}
  We say that the set \( \mscrU \) is a \term{Grothendieck universe} if it satisfied the following conditions:
  \begin{thmenum}
    \thmitem[def:grothendieck_universe/inductive]{GU1} It contains an \hyperref[def:inductive_set]{inductive set}.

    \thmitem[def:grothendieck_universe/transitive]{GU2} It is a \hyperref[def:transitive_set]{transitive set}.

    \thmitem[def:grothendieck_universe/power_set]{GU3} For any \( A \in \mscrU \), the \hyperref[def:basic_set_operations/power_set]{power set} \( \pow(A) \) also belongs to \( \mscrU \).

    \thmitem[def:grothendieck_universe/union]{GU4} For any member \( A \in \mscrU \) and any \( A \)-indexed family \( \set{ B_a }_{a \in A} \subseteq \mscrU \), the union
    \begin{equation*}
      \bigcup\set{ B_a \given a \in A }
    \end{equation*}
    belongs to \( \mscrU \). This is a restriction from unions over completely arbitrary families of sets to those families that can be indexed by members of \( A \).
  \end{thmenum}

  We formalize the entire concept via the following monstrous formula:
  \small
  \begin{equation*}\taglabel[\op{IsUniverse}]{eq:def:grothendieck_universe/predicate}
    \begin{aligned}
      \ref{eq:def:grothendieck_universe/predicate}[\upsilon] \coloneqq
        &\synneg \ref{eq:def:empty_set/predicate}[\upsilon]
        \synwedge
        \qforall {\tau \in \upsilon}
        \vast(
          \ref{eq:def:subset/predicate}[\tau, \upsilon]
          \synwedge
          \parens[\big]
            {
              \qexists {\synx \in \upsilon}
              \ref{eq:def:basic_set_operations/power_set/predicate}[\synx, \tau]
            }
          \synwedge \\ &\synwedge
          \qforall \synx
          \qforall \syny
          \parens[\Bigg]
          {
            \parens[\big]
              {
                \underbrace
                  {
                    \ref{eq:def:indexed_family/predicate}[\synx] \wedge \ref{eq:def:indexed_family/predicate_dom}[\synx, \tau] \wedge \ref{eq:def:indexed_family/predicate_ran}[\synx, \upsilon]
                  }_{\mathclap{\synx \T*{is a function} \tau \to \upsilon}}
                \synwedge
                \underbrace
                  {
                    \ref{eq:def:grothendieck_universe/predicate_isimage}[\syny, \synx]
                  }_{ \mathclap{\syny \T*{is the image of} \synx} }
              }
            \synimplies
            \underbrace
              {
                \qexists {\synz \in \upsilon} \ref{eq:def:basic_set_operations/union/predicate}[\synz, \synx]
              }_{ \bigcup \img(\synx) \in \upsilon}
        }
      \vast),
    \end{aligned}
  \end{equation*}
  \normalsize
  where
  \begin{equation*}\taglabel[\op{IsImage}]{eq:def:grothendieck_universe/predicate_isimage}
    \ref{eq:def:grothendieck_universe/predicate_isimage}[\rho, \tau]
    \coloneqq
    \qforall \synx
    \parens[\big]
    {
      \synx \in \rho
      \syniff
      \underbrace
      {
        \qexists {\syny \in \tau}
        \qexists \synz
        \ref{eq:def:kuratowski_pair/predicate}[\syny, \synz, \synx]
      }_{\tau(\synz) = \synx \T*{for some} \synx \in \dom(\tau) }
    }.
  \end{equation*}
\end{definition}

\begin{proposition}\label{thm:grothendieck_universe_closure}
  A \hyperref[def:grothendieck_universe]{Grothendieck universe} \( \mscrU \) has the following closure properties:
  \begin{thmenum}
    \thmitem{thm:grothendieck_universe_closure/subset} If \( A \) belongs to \( \mscrU \) and \( B \subseteq A \), the \( B \) also belongs to \( \mscrU \).

    \thmitem{thm:grothendieck_universe_closure/hereditarily_finite} Every \hyperref[def:universe_of_hereditary_finite_sets]{hereditarily finite set} belongs to \( \mscrU \).

    \thmitem{thm:grothendieck_universe_closure/binary_union} If \( A \) and \( B \) belong to \( \mscrU \), so does their union \( A \cup B \).

    \thmitem{thm:grothendieck_universe_closure/kuratowski_pair} If \( A \) and \( B \) belong to \( \mscrU \), so does the \hyperref[def:kuratowski_pair]{Kuratowski pair}
    \begin{equation*}
      \braket{ A, B } = \set{ \set{ A }, \set{ A, B } }.
    \end{equation*}

    \thmitem{thm:grothendieck_universe_closure/unordered_pair} If \( A \) and \( B \) belong to \( \mscrU \), so does the unordered pair \( \set{ A, B } \).

    \thmitem{thm:grothendieck_universe_closure/cartesian_product} If, for some \( \mscrK \)-\hyperref[def:indexed_family]{indexed family} \( \seq{ A_k }_{k \in \mscrK} \), both \( \mscrK \) and \( A_k \) belong to \( \mscrU \) for every \( k \in \mscrK \), the \hyperref[def:cartesian_product]{Cartesian product} \( \prod_{k \in \mscrK} A_k \) belongs to \( \mscrU \).

    \thmitem{thm:grothendieck_universe_closure/indexed_family} If, for some \( \mscrK \)-\hyperref[def:indexed_family]{indexed family} \( \seq{ A_k }_{k \in \mscrK} \), both \( \mscrK \) and \( A_k \) belong to \( \mscrU \) for every \( k \in \mscrK \), the \hyperref[def:indexed_family]{indexed family} \( \seq{ A_k }_{k \in \mscrK} \) belongs to \( \mscrU \).

    \thmitem{thm:grothendieck_universe_closure/function} If \( f: A \to B \) and both \( A \) and \( B \) belong to \( \mscrU \), so does \( f \).
  \end{thmenum}
\end{proposition}
\begin{proof}
  \SubProofOf{thm:grothendieck_universe_closure/subset} Suppose that \( B \subseteq A \) and that \( A \) is in \( \mscrU \).

  By the power set axiom \ref{def:grothendieck_universe/power_set}, \( \pow(A) \) belongs to \( \mscrU \). Since \( B \) in turn belongs to \( \pow(A) \), by the transitivity axiom \ref{def:grothendieck_universe/transitive}, it also belongs to \( \mscrU \).

  \SubProofOf{thm:grothendieck_universe_closure/hereditarily_finite} By axiom \ref{def:grothendieck_universe/inductive}, \( \mscrU \) contains \( \omega \).

  Since \( \omega \) contains \( 0 \), by the transitivity axiom \ref{def:grothendieck_universe/transitive}, \( 0 \) belongs to \( \mscrU \).

  With \( V_0 = \varnothing \) as a base case, via \fullref{thm:bounded_transfinite_induction} on \( n \in \omega \), from the power set axiom \ref{def:grothendieck_universe/power_set} it follows that \( V_{n+1} = \pow(V_n) \) is a member of \( \mscrU \) for every natural number \( n \).

  \SubProofOf{thm:grothendieck_universe_closure/binary_union} \Cref{thm:grothendieck_universe_closure/hereditarily_finite} implies that \( \set{ 0, 1 } \) belongs to \( \mscrU \).

  Via \ref{def:grothendieck_universe/union}, we conclude that the union \( A \cup B \) of the indexed family \( \set{ (0, A), (1, B) } \) also belongs to \( \mscrU \).

  \SubProofOf{thm:grothendieck_universe_closure/kuratowski_pair} \Cref{thm:grothendieck_universe_closure/binary_union} implies that \( A \cup B \) belongs to \( \mscrU \).

  Both \( A \) and \( B \) are subsets of \( A \cup B \), and hence members of \( \pow(A \cup B) \). The latter belongs to \( \mscrU \) due to the power set axiom \ref{def:grothendieck_universe/power_set}.

  Similarly, \( \set{ A } \) and \( \set{ A, B } \) are subsets of \( \pow(A \cup B) \), so they belong to \( \pow^2(A \cup B) \).

  Finally, the pair
  \begin{equation*}
    \braket{ A, B } = \set{ \set{ A }, \set{ A, B } }
  \end{equation*}
  belongs to \( \pow^3(A \cup B) \). Since the latter belongs to \( \mscrU \), by the transitivity axiom \ref{def:grothendieck_universe/transitive}, so does the pair itself.

  \SubProofOf{thm:grothendieck_universe_closure/unordered_pair} \Cref{thm:grothendieck_universe_closure/kuratowski_pair} implies that the Kuratowski pair \( \braket{ A, B } \) belongs to \( \mscrU \).

  Since \( \set{ A, B } \) is a subset of \( \braket{ A, B } \), by the transitivity axiom \ref{def:grothendieck_universe/transitive}, the unordered pair \( \set{ A, B } \) also belongs to \( \mscrU \).

  \SubProofOf{thm:grothendieck_universe_closure/cartesian_product} By the union schema \ref{def:grothendieck_universe/union}, the set
  \begin{equation*}
    \mscrA \coloneqq \bigcup_{k \in \mscrK} A_k
  \end{equation*}
  belongs to \( \mscrU \).

  \Cref{thm:grothendieck_universe_closure/binary_union} then implies that the binary union \( \mscrK \cup \mscrA \) belongs to \( \mscrU \). By using the power set axiom \ref{def:grothendieck_universe/power_set} several times, we conclude that \( \pow^4(\mscrK \cup \mscrA) \) is also in \( \mscrU \).

  The latter contains the Cartesian product \( \prod_{k \in \mscrK} A_k \). Indeed, any element of the product is a set of Kuratowski pairs \( \braket{ k, a_k } \), where \( k \) is from \( \mscrK \) and \( a_k \) is from \( A_k \subseteq \mscrA \).

  Such a pair
  \begin{equation*}
    \braket{ k, a_k } = \set{ \set{ k }, \set{ k, a_k } }
  \end{equation*}
  is a set of subsets of \( \mscrK \cup \mscrA \), i.e. an element of \( \pow^2(\mscrK \cup \mscrA) \).

  A single indexed family, i.e. a set of such pairs, belongs to \( \pow^3(\mscrK \cup \mscrA) \). The product, as a set of such indexed families, must thus belong to \( \pow^4(\mscrK \cup \mscrA) \).

  By the transitivity axiom \ref{def:grothendieck_universe/transitive}, the product \( \prod_{k \in \mscrK} A_k \) is in \( \mscrU \).

  \SubProofOf{thm:grothendieck_universe_closure/indexed_family} By \cref{thm:grothendieck_universe_closure/binary_union}, for every set \( A_k \), the singleton \( \set{ A_k } \) also belongs to \( \mscrU \).

  Then the Cartesian product \( \prod_{k \in \mscrK} \set{ A_k } \) contains exactly one member --- the desired indexed family \( \seq{ A_k }_{k \in K} \). \Cref{thm:grothendieck_universe_closure/cartesian_product} implies that the Cartesian product belongs to \( \mscrU \), hence, by the transitivity axiom \ref{def:grothendieck_universe/transitive}, so does the indexed family.

  \SubProofOf{thm:grothendieck_universe_closure/function} In \cref{def:set_valued_map}, we have formally defined a function \( f: A \to B \) as a triple \( (G, A, B) \), where \( G \) is a relation in \( A \times B \).

  \Cref{thm:grothendieck_universe_closure/cartesian_product} implies that \( A \times B \) is in \( \mscrU \), \cref{thm:grothendieck_universe_closure/subset} implies that its subset \( G \) is in \( \mscrU \), and finally \cref{thm:grothendieck_universe_closure/indexed_family} implies that the triple \( (G, A, B) \) is also in \( \mscrU \).
\end{proof}

\begin{proposition}\label{thm:smallest_grothendieck_universe_existence}
  Suppose that we are working in \logic{ZFC+U}. Then for any set \( A \), there exists a smallest Grothendieck universe containing \( A \).

  More generally, fix a set \( A \). Then there exists a smallest Grothendieck universe containing \( A \).
\end{proposition}
\begin{proof}
  If no set \( A \) is given, we simply take \( A = \varnothing \) since it must belong to every universe by definition.

  We use a trick analogous to \cref{thm:smallest_inductive_set_existence}.

  The \hyperref[def:zfcu]{axiom of universes} states that there exists at least one universe \( \mscrU \) that contains \( A \). Define
  \begin{equation*}
    \widehat \mscrU \coloneqq \set{ x \in \mscrU \given x \T{belongs to every Grothendieck universe} }.
  \end{equation*}

  Now that we have defined \( \widehat \mscrU \), it remains to verify that it is itself a universe. To show \ref{def:grothendieck_universe/inductive}, note that \( V_\omega \in \mscrU \) by \cref{thm:grothendieck_universe_closure/hereditarily_finite} and hence, by \ref{def:grothendieck_universe/transitive}, \( \omega \in \mscrU \).

  The rest of the verification is trivial.
\end{proof}

\begin{definition}\label{def:zfcu}\mcite[24]{MacLane1998CategoryTheory}
  The \term{axiom of universes} states that any set is contained in a \hyperref[def:grothendieck_universe]{Grothendieck universe}. Symbolically,
  \begin{equation}\label{eq:def:zfcu}
    \begin{aligned}
      \qforall \synx \qexists \upsilon \parens[\big]{ \ref{eq:def:grothendieck_universe/predicate}[\upsilon] \synwedge \synx \in \upsilon }.
    \end{aligned}
  \end{equation}

  We usually add this theorem to \hyperref[def:zfc]{ZFC} and call the resulting \hyperref[def:fol_theory]{logical theory} \logic{ZFC+U}.
\end{definition}

\begin{example}\label{ex:def:zfcu}
  From \cref{thm:grothendieck_universe_iff_regular_strong_limit} it follows that the universe of hereditarily finite sets \hyperref[def:universe_of_hereditary_finite_sets]{\( V_\omega \)} is a Grothendieck universe.

  The existence of other universes cannot be proven in \logic{ZFC}. For this reason, we use the \hyperref[def:zfcu]{axiom of universes}.
\end{example}

\begin{definition}\label{def:smallest_grothendieck_universe}\mimprovised
  In a \hyperref[def:fol_semantics/model]{model} of \hyperref[def:zfcu]{\logic{ZFC+U}}, we only know that at least one \hyperref[def:grothendieck_universe]{Grothendieck universe} \( \mscrU \) is required to exist.

  For definiteness, we can do the same trick as with \hyperref[def:inductive_set]{inductive sets}, taking the intersection of \( \mscrU \) with all other universes. The result is again a Grothendieck universe, and clearly the smallest. We denote it by lowercase \( \mscrU_0 \).
\end{definition}
\begin{comments}
  \item We use the smallest universe to determine whether a set is considered small in category theory --- see \cref{def:small_set}.
\end{comments}

\begin{proposition}\label{thm:grothendieck_universe_iff_regular_strong_limit}\mcite{Kruse1966GrothendieckUniverses}
  The stage \( V_\kappa \) of the \hyperref[def:cumulative_hierarchy]{von Neumann's cumulative hierarchy} is a \hyperref[def:grothendieck_universe]{Grothendieck universe} for every \hyperref[rem:strongly_inaccessible_cardinal]{regular strong limit cardinal} \( \kappa \).

  Conversely, for every Grothendieck universe \( \mscrU \), there exists a regular strong limit cardinal \( \kappa \) such that \( \mscrU = V_\kappa \).
\end{proposition}
\begin{proof}
  \SufficiencySubProof Let \( \kappa \) be a regular strong limit cardinal.

  \SubProofOf*{def:grothendieck_universe/inductive} Since \( \kappa \) is infinite and thus \( 0 < \kappa \), from \cref{thm:def:cumulative_hierarchy/membership} it follows that \( V_0 \in V_\kappa \).

  \SubProofOf*{def:grothendieck_universe/transitive} The set \( V_\kappa \) is transitive as shown in \cref{thm:def:cumulative_hierarchy/transitive}.

  \SubProofOf*{def:grothendieck_universe/power_set} Since \( \kappa \) is a limit ordinal, \( V_\kappa \) satisfies the axiom of power sets as shown in \fullref{thm:cumulative_hierarchy_model_of_z} and thus if \( A \in V_\kappa \), then \( \pow(A) \in V_\kappa \).

  \SubProofOf*{def:grothendieck_universe/union} Fix some member \( A \in V_\kappa \) and some \( A \)-indexed family \( \set{ B_a }_{a \in A} \subseteq V_\kappa \). From \cref{thm:strong_regular_cardinal_stage_cardinality} it follows that \( \card(A) < \kappa \) and \( \card(B_a) < \kappa \) for every \( a \in A \). Thus,
  \begin{equation*}
    \card(\set{ B_a \given a \in A }) \leq \card(A) < \kappa
  \end{equation*}
  and from \cref{thm:regular_cardinal_stage_inverse_transitivity} it follows that
  \begin{equation*}
    \set{ B_a \given a \in A } \in V_\kappa.
  \end{equation*}

  The union
  \begin{equation*}
    \bigcap \set{ B_a \given a \in A }
  \end{equation*}
  is then a member of a lower stage, hence it also belongs to \( V_\kappa \).

  \NecessitySubProof Let \( \mscrU \) be a Grothendieck universe and let \( \alpha \) be the smallest \emph{ordinal} not in \( \mscrU \). We will first show that \( \mscrU = V_\alpha \) and gradually prove that \( \alpha \) is actually an inaccessible cardinal.

  \SubProof*{Proof that \( V_\beta \in \mscrU \) for ordinals \( \beta \in \mscrU \)} We will use \fullref{thm:bounded_transfinite_induction} on \( \beta < \alpha \).
  \begin{itemize}
    \item From \cref{thm:grothendieck_universe_closure/hereditarily_finite} it follows that \( \varnothing \in \mscrU \).

    \item If \( \beta < \alpha \) and \( V_\beta \in \mscrU \), then by \ref{def:grothendieck_universe/power_set} we have \( V_{\beta + 1} = \pow(V_\beta) \in \mscrU \).

    We will come back to this step a bit later, but for now note that \( V_{\beta + 1} \in \mscrU \) regardless of whether \( \beta + 1 \in \mscrU \).

    \item If \( \lambda < \alpha \) is a limit ordinal and \( V_\beta \in \mscrU \) for every \( \beta < \lambda \), we have
    \begin{equation*}
      V_\lambda
      \reloset {\eqref{eq:def:cumulative_hierarchy}} =
      \bigcup\set{ V_\beta \given \beta < \lambda },
    \end{equation*}
    which is a \( \lambda \)-indexed union of members of \( \mscrU \). Since \( \lambda \in \mscrU \), by \ref{def:grothendieck_universe/power_set} we have \( V_\lambda \in \mscrU \).
  \end{itemize}

  \SubProof*{Proof that \( \alpha \) is a limit ordinal} In the successor case we noted that \( V_{\beta + 1} \in \mscrU \) for every \( \beta < \alpha \) regardless of whether \( \beta + 1 < \alpha \). Since \( \rank(\beta + 1) = \beta + 1 \), it follows that \( \beta + 1 \in V_{\beta + 2} \in \mscrU \) and thus by \ref{def:grothendieck_universe/transitive}, \( \beta + 1 \in \mscrU \). Therefore, \( \alpha \) cannot be a successor ordinal --- if \( \alpha = \beta + 1 \), then \( \beta \in \mscrU \) by definition of \( \alpha \) and thus \( \beta + 1 = \alpha \in \mscrU \), which is a contradiction.

  Since \( \alpha > 0 \), it remains for \( \alpha \) to be a limit ordinal.

  \SubProof*{Proof that \( V_\alpha \subseteq \mscrU \)} By \ref{def:grothendieck_universe/transitive}, \( V_\beta \subseteq \mscrU \) for every \( \beta < \alpha \). We can conclude that
  \begin{equation*}
    V_\alpha
    \reloset {\eqref{eq:def:cumulative_hierarchy}} =
    \bigcup\set{ V_\beta \given \beta < \alpha }
    \subseteq
    U.
  \end{equation*}

  In order to show that equality holds, we must first prove that \( \alpha \) is a strongly inaccessible cardinal. But this requires some auxiliary results.

  \SubProof*{Proof that \( \set{ B } \in \mscrU \) for every \( B \in \mscrU \)} By \ref{def:grothendieck_universe/power_set} we have that \( \pow(\pow(B)) \in \mscrU \). But \( \set{ B } \subseteq \pow(B) \) and hence \( \set{ B } \in \pow(\pow(B)) \). By \ref{def:grothendieck_universe/transitive}, \( \set{ B } \in \mscrU \).

  \SubProof*{Proof that \( \kappa = \alpha \) is a cardinal} Suppose that \( \alpha \) is not a cardinal. Indeed, suppose that there exists some \( \beta < \alpha \) such that there exists a bijective function \( f: \beta \to \alpha \). Then
  \begin{equation*}
    \alpha = \bigcup\set{ \set{ f(\gamma) } \given \gamma < \beta }
  \end{equation*}
  is a \( \beta \)-indexed union of members of \( \alpha \) and hence \( \alpha \in \alpha \). But this contradicts \cref{thm:simple_foundation_theorems/member_of_itself}.

  Therefore, \( \alpha \) is a cardinal. We will henceforth denote it by \( \kappa \) to highlight that it is a cardinal.

  \SubProof*{Proof that \( \card(B) \in \mscrU \) for every \( B \in \mscrU \)} Let \( B \in \mscrU \) and let \( f: B \to \card(B) \) be a bijective function. Then
  \begin{equation*}
    \card(B)
    =
    f[B]
    =
    \set{ f(x) \given x \in B }
    =
    \bigcup\set[\big]{ \set{ f(x) } \given x \in B }
  \end{equation*}
   is a \( B \)-indexed union of members of \( \mscrU \) and, by \ref{def:grothendieck_universe/union}, \( \card(B) \in \mscrU \).

  \SubProof*{Proof that \( \kappa \) is a strong limit} For every \( \beta < \kappa \) by \ref{def:grothendieck_universe/power_set} we have \( \pow(\beta) \in \mscrU \). We have already shown that \( \card(\pow(\beta)) \in \mscrU \) and, by \cref{thm:cardinal_exponentiation_power_set}, we have
  \begin{equation*}
    \card(\pow(\beta)) = 2^{\card(\beta)} = 2^\beta.
  \end{equation*}

  Hence, \( 2^\beta < \kappa \) and \( \kappa \) is a strong limit.

  \SubProof*{Proof that \( \kappa \) is regular} Let \( C \subseteq \kappa \) be an unbounded set. We will show that \( \card(C) = \kappa \).

  Suppose that \( \card(C) < \kappa \). Then \( \card(C) \in \mscrU \) since \( \kappa = \alpha \) is the smallest ordinal not contained in \( \mscrU \). Let \( f: \card(C) \to C \) be a bijective function. Then
  \begin{equation*}
    C = \bigcup\set[\big]{ \set{ f(\gamma) } \given \gamma < \card(C) }
  \end{equation*}
  and by \ref{def:grothendieck_universe/union}, \( C \in \mscrU \).

  Since \( C \) is unbounded, we have \( \sup C \geq \kappa \). But from \cref{thm:union_of_set_of_ordinals} is follows that \( \bigcup C = \sup C \) and by \ref{def:grothendieck_universe/union}, \( \bigcup C \in \mscrU \), which implies that \( \sup C < \kappa \).

  The obtained contradiction shows that \( \card(C) = \kappa \). Since \( C \) was an arbitrary unbounded set, it follows that \( \kappa \) satisfies \cref{def:regular_cardinal/unbounded_subsets} and is thus regular.

  \SubProof*{Proof that \( V_\kappa = U \)} Finally, now that we know that \( \kappa \) is a strongly inaccessible cardinal, we can show that equality holds in \( V_\kappa \subseteq \mscrU \).

  Aiming at a contradiction, suppose that \( U \setminus V_\kappa \) is nonempty. By the \hyperref[def:zfc/foundation]{axiom of foundation}, there exists a set \( C \in \mscrU \setminus V_\kappa \) such that
  \begin{equation*}
    C \cap (U \setminus V_\kappa) = \varnothing,
  \end{equation*}
  thus \( C \subseteq V_\kappa \). From \cref{thm:strong_regular_cardinal_stage_cardinality} if follows that \( \card(C) < \kappa \) and from \cref{thm:regular_cardinal_stage_inverse_transitivity} it follows that \( C \in V_\kappa \), which contradicts our choice of \( C \) as a member of \( U \setminus V_\kappa \).

  Therefore, \( V_\kappa = U \).
\end{proof}

\begin{corollary}\label{thm:grothendieck_universe_is_model_of_zfc}
  Every uncountable Grothendieck universe is a standard model of \hyperref[def:zfc]{\logic{ZFC}}.
\end{corollary}
\begin{proof}
  Follows from \cref{thm:grothendieck_universe_iff_regular_strong_limit} and \fullref{thm:cumulative_hierarchy_model_of_zfc}.
\end{proof}
